\documentclass[11pt]{article}

\usepackage[utf8]{inputenc}
\usepackage[T1]{fontenc}
\usepackage{amsmath,amssymb,amsthm,mathtools}
\usepackage{enumitem}

\newtheorem{theorem}{Theorem}
\newtheorem{lemma}[theorem]{Lemma}
\newtheorem{proposition}[theorem]{Proposition}
\newtheorem{definition}[theorem]{Definition}
\newtheorem{remark}[theorem]{Remark}

\newcommand{\R}{\mathbb R}
\newcommand{\Z}{\mathbb Z}
\newcommand{\len}{\operatorname{length}}


\begin{document}
\title{Submission C solution to Problem 2}
\date{}
\maketitle

\begin{abstract}
We prove the stronger statement that a number \(\beta>0\) is realized in the sense of the question if and only if \(\beta>\sqrt 3\).  Hence the infimum of the set of realized numbers is \(\sqrt 3\).
\end{abstract}

\section{Preliminaries}

Let
\[
M_\beta=\Sigma/\langle G_\beta\rangle .
\]
Thus \(M_\beta\) is the flat Mobius band obtained from the rectangle
\([0,\beta]\times[0,1]\) by identifying \((0,y)\sim(\beta,1-y)\).  Its boundary is a circle of flat length \(2\beta\).  We use the boundary coordinate
\[
s(x,0)=x\pmod {2\beta},\qquad s(x,1)=x+\beta\pmod {2\beta}.
\]

\begin{lemma}[quotient and boundary contraction]
Suppose \(\beta\) is realized.  Then the map \(f\) descends to a topological embedding
\[
\bar f:M_\beta\longrightarrow \R^3 .
\]
Moreover, the \(G_\beta\)-invariant clean triangulation descends to a finite clean triangulation of \(M_\beta\), and every nondegenerate boundary arc \(A\subset \partial M_\beta\) satisfies
\[
\len(\bar f(A))<\len_0(A),
\]
where \(\len_0\) denotes flat boundary length.
\end{lemma}

\begin{proof}
Condition (2) gives \(f\circ G_\beta=f\), so \(f\) descends to \(\bar f\).  The same condition says precisely that \(\bar f\) is injective.  Since \(M_\beta\) is compact and \(\R^3\) is Hausdorff, \(\bar f\) is an embedding.

A compact fundamental domain in \(\Sigma\) meets only finitely many triangles by local finiteness, so only finitely many \(G_\beta\)-orbits of triangles occur.  Thus the quotient triangulation is finite.

Finally, on each boundary edge of the triangulation, \(f\) is affine and strictly shortens that edge.  Hence it strictly shortens every subsegment of every boundary edge.  Since every boundary arc is a finite union of such subsegments, its image has strictly smaller length.
\end{proof}

\begin{lemma}[the squeeze cone]
Let \(\tau\) be a clean triangle with boundary edge \(AB\) and opposite vertex \(C\).  If an affine map \(L:\tau\to \R^3\) decreases the length of \(AB\) and increases the lengths of \(AC\) and \(BC\), then \(L\) increases the length of every segment joining \(C\) to a point of \(AB\).
\end{lemma}

\begin{proof}
Write \(u=C-A\), \(v=C-B\), so \(u-v=B-A\).  Let
\[
q(w)=|Lw|^2-|w|^2 .
\]
The hypotheses are
\[
q(u)>0,\qquad q(v)>0,\qquad q(u-v)<0.
\]
For \(0\le t\le 1\),
\[
q((1-t)u+tv)
=(1-t)q(u)+tq(v)-t(1-t)q(u-v)>0.
\]
Thus every cevian from \(C\) to the base \(AB\) is strictly lengthened.
\end{proof}

\section{The canonical ruling and a \(T\)-pattern}

For every clean triangle, foliate it by the straight segments joining the vertex opposite the boundary edge to points of the boundary edge.  Adjacent clean triangles share a nonboundary edge, and that edge is one of the leaves in both triangles.  Therefore these local foliations glue.

\begin{lemma}[canonical ruling]
The quotient clean triangulation of \(M_\beta\) carries a continuous circle-family \(\mathcal L=\{L_t\}_{t\in S^1}\) of straight clean segments with the following properties.
\begin{enumerate}[label=\textup{(\roman*)}]
\item Each \(L_t\) is contained in a single clean triangle and joins the two boundary components in the universal cover.
\item The relative interiors of distinct leaves are disjoint and the leaves cover \(M_\beta\).
\item The image \(\bar f(L_t)\) is a straight segment in \(\R^3\), and
\[
|\bar f(L_t)|>|L_t|_0\ge 1.
\]
\item A continuous orientation of the leaves reverses after one circuit of the parameter circle.
\end{enumerate}
\end{lemma}

\begin{proof}
Let \(\mathcal T\) be the finite clean triangulation of \(M_\beta\).  Form the dual graph whose vertices are triangles and whose edges correspond to shared nonboundary edges.  Every triangle has exactly two nonboundary edges, and every nonboundary edge is shared by two triangles, so the dual graph is \(2\)-regular.  It is connected: at each boundary vertex the link is an interval, so the triangles incident to that vertex are connected through shared nonboundary edges; since the triangulation is connected, the dual graph is connected.  Hence the dual graph is a single cycle.

On a triangle \(\tau\), the family of cevians from the opposite vertex to the boundary edge is parametrized by an interval; its two endpoint leaves are the two nonboundary edges of \(\tau\).  Gluing these intervals according to the cyclic dual graph gives a parameter circle \(S^1\).  This proves the existence and continuity of the family \(\{L_t\}\), and the relative interiors are plainly disjoint.

Each leaf is contained in a clean triangle, and by the preceding squeeze-cone lemma its image is strictly longer than its flat length.  Since every clean segment crosses from \(y=0\) to \(y=1\) in the universal cover, its flat length is at least \(1\).

Finally, the leaf family is an interval-bundle structure on the Mobius band away from finitely many boundary vertices.  Since \(M_\beta\) is the nontrivial \(I\)-bundle over a circle, the orientation line of the fibers is nontrivial.  Equivalently, in the universal cover, going once around the leaf circle applies the deck transformation \(G_\beta\), which exchanges the two boundary components.  Hence any continuous orientation of the leaves is reversed after one circuit.
\end{proof}

\begin{lemma}[\(T\)-pattern]
There are two distinct leaves \(L_s,L_t\) such that the two straight segments
\[
\bar f(L_s),\qquad \bar f(L_t)
\]
lie in perpendicular coplanar lines.
\end{lemma}

\begin{proof}
Lift the parameter circle to \(\R\).  Choose continuously oriented unit direction vectors \(u(r)\) for the image segments and let \(m(r)\) denote their midpoints.  By the orientation-reversal property,
\[
u(r+2\pi)=-u(r),\qquad m(r+2\pi)=m(r).
\]

Let
\[
\Upsilon=\{(r,s): r<s<r+2\pi\}/(r,s)\sim(r+2\pi,s+2\pi).
\]
Compactify \(\Upsilon\) by adding two points: \(\partial_+\), corresponding to \(s\to r\), and \(\partial_-\), corresponding to \(s\to r+2\pi\).  This compactification is a \(2\)-sphere, and the swap involution
\[
(r,s)\longmapsto (s,r+2\pi)
\]
extends to the antipodal map on that sphere.

Define
\[
F(r,s)=\left(u(r)\cdot u(s),\ (m(r)-m(s))\cdot (u(r)\times u(s))\right).
\]
Then \(F\) extends continuously to the compactification with
\[
F(\partial_+)=(1,0),\qquad F(\partial_-)=(-1,0),
\]
and satisfies
\[
F(s,r+2\pi)=-F(r,s).
\]
By the Borsuk--Ulam theorem, \(F\) has a zero.  At such a zero, the two direction vectors are perpendicular, and the vector joining the midpoints is orthogonal to the normal \(u(r)\times u(s)\).  Therefore the two supporting lines are coplanar.  This gives the desired pair of leaves.
\end{proof}

\section{The lower bound}

\begin{lemma}[geometry of a squeezed \(T\)-pattern]
If \(\beta\) is realized, then \(\beta>\sqrt 3\).
\end{lemma}

\begin{proof}
Let \(T\) and \(B\) be the two leaves supplied by the preceding lemma.  Write
\[
T^*=\bar f(T),\qquad B^*=\bar f(B).
\]
Their supporting lines are perpendicular and coplanar.  Interchange \(T\) and \(B\), if necessary, so that the line supporting \(T^*\) does not meet the relative interior of \(B^*\).  Otherwise the two embedded segments would have an interior intersection.

Let \(u\) be the endpoint of \(B\) whose image is farthest from the line supporting \(T^*\).  Since \(B^*\) is perpendicular to that line and the line does not meet the relative interior of \(B^*\), the distance \(h\) from \(\bar f(u)\) to the line supporting \(T^*\) satisfies
\[
h\ge |B^*|>|B|_0\ge 1.
\]

Let the endpoints of \(T\) divide \(\partial M_\beta\) into two boundary arcs \(H\) and \(D\), with \(u\in D\).  Choose the signed horizontal displacement \(t\) of \(T\) so that
\[
\len_0(H)=\beta-t,\qquad \len_0(D)=\beta+t.
\]
Then
\[
|T|_0=\sqrt{1+t^2}.
\]
Because \(T\) is a leaf, it is lengthened:
\[
|T^*|>\sqrt{1+t^2}.
\]
Since \(H\) is a boundary path joining the endpoints of \(T\),
\[
\sqrt{1+t^2}<|T^*|\le \len(\bar f(H))<\beta-t. \tag{1}
\]

Now consider the planar triangle \(\Delta\) whose base is \(T^*\) and whose third vertex is \(\bar f(u)\).  Its base has length \(>|T|_0=\sqrt{1+t^2}\), and its height is \(h>1\).  If \(\vee\) denotes the union of the two nonbase sides of \(\Delta\), the elementary reflection argument gives
\[
\len(\vee)>\sqrt{5+t^2}.
\]
Indeed, among all triangles with base length \(b\) and height \(h\), the sum of the two nonbase sides is at least \(\sqrt{b^2+4h^2}\).

The boundary arc \(D\) runs from one endpoint of \(T\) to the other through \(u\).  Hence the two sides of \(\vee\) are no longer than the two corresponding subarcs of \(\bar f(D)\).  Therefore
\[
\sqrt{5+t^2}<\len(\vee)\le \len(\bar f(D))<\beta+t. \tag{2}
\]
From (1) and (2),
\[
2\beta>\alpha(t):=\sqrt{1+t^2}+\sqrt{5+t^2},
\]
and also
\[
2\beta>\gamma(t):=2\sqrt{5+t^2}-2t.
\]
Let \(t_0=1/\sqrt3\).  Then
\[
\alpha(t_0)=\gamma(t_0)=2\sqrt3.
\]
Moreover, \(\alpha\) is increasing for \(t>0\), and \(\gamma\) is strictly decreasing on \(\R\).  Thus for every real \(t\),
\[
\max\{\alpha(t),\gamma(t)\}\ge 2\sqrt3.
\]
Consequently \(2\beta>2\sqrt3\), so \(\beta>\sqrt3\).
\end{proof}

\section{Realization for every \(\beta>\sqrt3\)}

We use the classical existence theorem for paper Mobius bands: for every width \(w>0\) and every length \(\lambda\) with \(\lambda/w>\sqrt3\), there exists a smooth embedded isometric Mobius band of width \(w\) and length \(\lambda\), obtained by smoothing the triangular Mobius band.  Moreover, it has a continuous ruling by straight bends whose preimages are straight segments joining the two boundary components.  This is the standard Sadowsky--Wunderlich construction; see Sadowsky \cite{Sadowsky}, the translation by Hinz--Fried \cite{HinzFried}, and the modern discussion in Schwartz \cite{Schwartz}.

\begin{proposition}
Every \(\beta>\sqrt3\) is realized.
\end{proposition}

\begin{proof}
Fix \(\beta>\sqrt3\).  Choose numbers \(w,\lambda\) such that
\[
1<w,\qquad \sqrt3\,w<\lambda<\beta,\qquad
w^2-1>\beta^2-\lambda^2.
\]
This is possible by taking \(\lambda\) sufficiently close to \(\beta\) and then \(w>1\) sufficiently close to \(1\).

Let \(\rho=\lambda/\beta<1\).  Let
\[
P_{\lambda,w}=([0,\lambda]\times[0,w])/(0,y)\sim(\lambda,w-y),
\]
and let
\[
I:P_{\lambda,w}\longrightarrow \R^3
\]
be a smooth embedded paper Mobius band with a ruling by straight bends.  Define
\[
A:M_\beta\longrightarrow P_{\lambda,w},\qquad A[x,y]=[\rho x,wy].
\]
This is well-defined because
\[
A(x+\beta,1-y)=(\rho x+\lambda,w-wy),
\]
which is the defining Mobius identification in \(P_{\lambda,w}\).  Put
\[
F=I\circ A:M_\beta\longrightarrow \R^3.
\]
Then \(F\) is an embedding.

Let \(\{C_s\}_{s\in S^1}\) be the pulled-back bend ruling on \(M_\beta\).  Each \(C_s\) is a clean straight segment.  If \(C_s\) has signed horizontal displacement \(d\), with \(|d|\le \beta\), then
\[
|F(C_s)|^2=\rho^2d^2+w^2.
\]
Hence
\[
|F(C_s)|^2-|C_s|_0^2
=\rho^2d^2+w^2-(d^2+1)
=(w^2-1)-(1-\rho^2)d^2.
\]
Since \(|d|\le\beta\),
\[
(1-\rho^2)d^2\le (1-\rho^2)\beta^2=\beta^2-\lambda^2,
\]
so
\[
|F(C_s)|^2-|C_s|_0^2\ge w^2-1-(\beta^2-\lambda^2)>0. \tag{3}
\]
Thus every pulled-back bend is strictly lengthened by \(F\).

Choose finitely many bends \(C_{s_0},\dots,C_{s_{N-1}}\) in cyclic order, with mesh sufficiently small.  Consecutive bends bound clean quadrilaterals in \(M_\beta\).  Split each such quadrilateral by one of its diagonals.  This gives a finite clean triangulation of \(M_\beta\).  Define a piecewise-affine map \(\bar f\) by setting \(\bar f=F\) on all vertices and extending affinely over each triangle.

For sufficiently small mesh, \(\bar f\) is an embedding.  Indeed, in bend coordinates the smooth map \(F\) has the form
\[
F(s,r)=(1-r)a(s)+r b(s),
\]
with \(a,b\) \(C^1\) boundary curves.  The above piecewise-affine maps are the usual secant approximations to this ruled embedding; they converge to \(F\) in \(C^1\).  A \(C^1\)-small perturbation of a compact \(C^1\)-embedded surface is again embedded.

Now check the squeeze inequalities.  If \(e\) is a boundary edge of flat length \(\ell\), then the boundary arc of \(F\) between its endpoints has length \(\rho\ell\), so the chord \(\bar f(e)\) has length at most \(\rho\ell<\ell\).  Hence boundary edges are strictly shortened.

The nonboundary edges are of two types: selected bends and diagonals between two sufficiently close selected bends.  Selected bends are strictly lengthened by (3).  For the diagonals, the relevant squared length difference depends continuously on the two adjacent bend parameters and agrees, when the parameters coincide, with the positive quantity in (3).  Compactness of the parameter circle gives a uniform positive margin, so for sufficiently small mesh every diagonal is also strictly lengthened.

Therefore \(\bar f\) restricts to a squeeze map on every triangle of the clean triangulation of \(M_\beta\).  Lifting the triangulation and the map to \(\Sigma\) gives a \(G_\beta\)-invariant clean triangulation and a continuous map
\[
f:\Sigma\to\R^3
\]
with \(f=\bar f\circ \pi\), where \(\pi:\Sigma\to M_\beta\) is the quotient map.  Since \(\bar f\) is an embedding,
\[
f(p)=f(q)\quad\Longleftrightarrow\quad \pi(p)=\pi(q)
\quad\Longleftrightarrow\quad p=G_\beta^k(q)\text{ for some }k\in\Z.
\]
Thus \(\beta\) is realized.
\end{proof}

\begin{theorem}
The set of realized numbers is exactly \((\sqrt3,\infty)\).  In particular, its infimum is \(\sqrt3\).
\end{theorem}

\begin{proof}
The lower-bound lemma shows that every realized \(\beta\) satisfies \(\beta>\sqrt3\).  The proposition shows that every \(\beta>\sqrt3\) is realized.  Hence the set of realized numbers is \((\sqrt3,\infty)\), whose infimum is \(\sqrt3\).
\end{proof}

\begin{thebibliography}{9}

\bibitem{Sadowsky}
M. Sadowsky,
\emph{Ein elementarer Beweis fuer die Existenz eines abwickelbaren Moebiusschen Bandes und die Zurueckfuehrung des geometrischen Problems auf ein Variationsproblem},
Sitzungsberichte der Preussischen Akademie der Wissenschaften, physikalisch-mathematische Klasse \textbf{22} (1930), 412--415.

\bibitem{HinzFried}
D. F. Hinz and E. Fried,
\emph{Translation of Michael Sadowsky's paper ``An elementary proof for the existence of a developable Moebius band and the attribution of the geometric problem to a variational problem''},
J. Elasticity \textbf{119} (2015), 3--6.

\bibitem{Schwartz}
R. E. Schwartz,
\emph{The optimal paper Moebius band},
Ann. of Math. (2) \textbf{201} (2025), no. 1, 291--305.

\end{thebibliography}

\end{document}